% Part: computability
% Chapter: computability-theory
% Section: equiv-ce-defs

\documentclass[../../../include/open-logic-section]{subfiles}

\begin{document}

\olfileid{cmp}{thy}{eqc}

\olsection[గణనీయంగా లెక్కించదగిన సమితుల నిర్వచనలు]{గణనీయంగా
  లెక్కించదగిన సమితుల తుల్య నిర్వచనలు}


గణనీయంగా లెక్కించదగినది కావడం అంటే ఏమిటో చెప్పే అనేక ముఖ్యమైన తుల్య
ప్రకటనలను కింద ఇస్తున్నాం.
\sourcecorrection{OLTECOMTHY-005}{ఆంగ్ల మూలంలోని పూర్తి విభాగ శీర్షికలో ``నిర్వచనలు'' అనే పదానికి అక్షరదోషం ఉంది. సంక్షిప్త శీర్షికకూ తరువాతి తుల్యతా సిద్ధాంతానికీ అనుగుణంగా సరైన రూపాన్ని అనువదించాం.}

\begin{thm}
\ollabel{thm:ce-equiv}
$S$ సహజ సంఖ్యల సమితి అనుకుందాం. అప్పుడు కింది ప్రకటనలు తుల్యాలు:
\begin{enumerate}
\item\ollabel{case:ce} $S$ గణనీయంగా లెక్కించదగినది.
\item\ollabel{case:ran-pc} $S$ ఒక \emph{పాక్షిక} గణనీయ ప్రమేయపు వ్యాప్తి.
\item\ollabel{case:ran-prim} $S$ శూన్య సమితి, లేదా ఒక ఆదిమ పునరావృత్త ప్రమేయపు వ్యాప్తి.
\item\ollabel{case:ce-domain} $S$ ఒక పాక్షిక గణనీయ ప్రమేయపు \emph{నిర్వచన క్షేత్రం}.
\end{enumerate}
\end{thm}

\begin{explain}
మొదటి మూడు ఉపవాక్యాలు, ఏ శూన్యేతర గణనీయంగా లెక్కించదగిన సమితినైనా
ఒక గణనీయ ప్రమేయం, ఒక పాక్షిక గణనీయ ప్రమేయం, లేదా ఒక ఆదిమ పునరావృత్త
ప్రమేయం ద్వారా లెక్కించవచ్చని తుల్యంగా తీసుకోవచ్చని చెబుతాయి. నాలుగవ
ఉపవాక్యం $S$ గణనీయంగా లెక్కించదగినదైతే, ఏదో సూచిక~$e$కు
\[
S=\Setabs{x}{\cfind{e}(x)\fdefined}.
\]
అవుతుందని చెబుతుంది. మరో మాటలో, $S$ అనేది $\cfind{e}$ గణన ఆగే
ఇన్‌పుట్‌ల సమితి. అందుకే గణనీయంగా లెక్కించదగిన సమితులను కొన్నిసార్లు
\emph{అర్ధ-నిర్ణయించదగినవి} అంటారు: ఒక సంఖ్య సమితిలో ఉంటే చివరకు
``అవును'' అనే సమాధానం వస్తుంది; లేకపోతే ``కాదు'' అనే సమాధానం ఎన్నటికీ
రాదు!{}
\end{explain}

\begin{proof}
ప్రతి ఆదిమ పునరావృత్త ప్రమేయం గణనీయం; ప్రతి గణనీయ ప్రమేయం పాక్షిక
గణనీయం. కనుక \olref{case:ran-prim} నుంచి \olref{case:ce},
\olref{case:ce} నుంచి~\olref{case:ran-pc} వస్తాయి. ($S$ శూన్య సమితి
అయితే, ఎక్కడా నిర్వచితం కాని పాక్షిక గణనీయ ప్రమేయపు వ్యాప్తి~$S$ అని
గమనించండి.) \olref{case:ran-pc} నుంచి \olref{case:ran-prim} వస్తుందని
చూపితే, మొదటి మూడు ఉపవాక్యాలు తుల్యాలని చూపినట్లే.

కాబట్టి $S$ అనేది పాక్షిక గణనీయ ప్రమేయం $\cfind{e}$ యొక్క వ్యాప్తి
అనుకుందాం. $S$ శూన్య సమితి అయితే పని పూర్తయింది. లేకపోతే $S$లో ఏదైనా
మూలకం $a$ను తీసుకోండి. క్లీనీ నియత రూప సిద్ధాంతం ప్రకారం,
\[
\cfind{e}(x)=U(\umin{s}{T(e,x,s)}).
\]
అని రాయవచ్చు. ముఖ్యంగా, $T(e,x,s)$ నెరవేరి $U(s)=y$ అయ్యే $s$
ఒకటి ఉంటే, అప్పుడే మాత్రమే $\cfind{e}(x)\fdefined$ అవుతుంది మరియు
దాని విలువ $=y$. $f(z)$ను ఇలా నిర్వచించండి:
\[
f(z)=\begin{cases}
  U((z)_1) & \text{$T(e,(z)_0,(z)_1)$ అయితే} \\
  a & \text{ఇతర సందర్భాల్లో.}
\end{cases}
\]
$T$, $U$లు ఆదిమ పునరావృత్తాలు కనుక $f$ కూడా ఆదిమ పునరావృత్తం.
\iftag{TMs}{$z$ అనేది $\tuple{(z)_0,(z)_1}$ జతను సంకేతీకరిస్తూ,
  $(z)_1$ అనేది ఇన్‌పుట్ $(z)_0$పై యంత్రం~$M_e$ యొక్క ఆగే గణన
  అయితే, ట్యూరింగ్ యంత్రాల పరంగా $f$ ఆ గణన అవుట్‌పుట్‌ను తిరిగి
  ఇస్తుంది; లేకపోతే $a$ను తిరిగి ఇస్తుంది.}

$S$ అనేది $f$ యొక్క వ్యాప్తి అని చూపాలి; అంటే ఏ సహజ సంఖ్య~$y$కైనా,
$y\in S$ అయితేనే అది $f$ వ్యాప్తిలో ఉంటుందని చూపాలి. ముందు దిశలో,
$y\in S$ అనుకుందాం. అప్పుడు $y$ అనేది $\cfind{e}$ వ్యాప్తిలో ఉంది;
కాబట్టి ఏవో $x$,~$s$లకు $T(e,x,s)$ నెరవేరి $U(s)=y$ అవుతుంది.
అప్పుడు $y=f(\tuple{x,s})$. విలోమంగా, $y$ అనేది $f$ వ్యాప్తిలో
ఉందనుకుందాం. అప్పుడు $y=a$, లేదా ఏదో~$z$కు $T(e,(z)_0,(z)_1)$
నెరవేరి $U((z)_1)=y$ అవుతుంది. రెండవ సందర్భంలో
$\cfind{e}((z)_0)\fdefined=y$; కనుక రెండు సందర్భాల్లోనూ $y$ అనేది~$S$లో ఉంది.
\sourcecorrection{OLTECOMTHY-006}{ఈ విలోమ దిశలో ఆంగ్ల మూలం $T(e,(z)_0,(z)_1)$ నుంచి వచ్చే నిర్వచిత విలువను $\cfind{e}(x)$గా రాసింది; ఆ పేరాలో $x$కు విలువ కట్టలేదు. చూపించిన సంకేతీకృత జతకు అనుగుణంగా దాన్ని $\cfind{e}((z)_0)$గా సరిచేశాం.}

($\cfind{e}(x)\fdefined=y$ అనే సంకేతం ``$\cfind{e}(x)$ నిర్వచితమై,
దాని విలువ $y$కు సమానం'' అని అర్థం. బదులుగా $\cfind{e}(x)=y$ అని
రాయవచ్చు; కానీ అదనపు బాణం మనం పాక్షిక ప్రమేయంతో వ్యవహరిస్తున్నామని
గుర్తుచేయడానికి కొన్నిసార్లు సహాయపడుతుంది.)

\olref{thm:ce-equiv} నిరూపణను పూర్తి చేయడానికి, \olref{case:ce},
\olref{case:ce-domain} తుల్యాలని చూపితే చాలు. మొదట \olref{case:ce}
నుంచి~\olref{case:ce-domain} వస్తుందని చూపుదాం. $S$ అనేది గణనీయ
ప్రమేయం~$f$ యొక్క వ్యాప్తి అనుకుందాం, అంటే
\[
S=\Setabs{y}{\text{ఏదో $x$కు, }f(x)=y}.
\]
ఇలా ఉంచండి:
\[
g(y)=\umin{x}{(f(x)=y)}.
\]
అప్పుడు $g$ పాక్షిక గణనీయ ప్రమేయం; ఏదో~$x$కు $f(x)=y$ అయితేనే
$g(y)$ నిర్వచితమవుతుంది. మరో మాటలో, $g$ నిర్వచన క్షేత్రమే $f$
వ్యాప్తి. \iftag{TMs}{ట్యూరింగ్ యంత్రాల పరంగా: $S$లోని మూలకాలను
లెక్కించే ట్యూరింగ్ యంత్రం~$F$ ఇచ్చినప్పుడు, ఇచ్చిన మూలకం సమితిలో
ఉందో లేదో తెలుసుకోవడానికి $F$ అవుట్‌పుట్‌లలో వెతుకుతూ, ఉంటే ఆగి,
లేకపోతే ఎప్పటికీ వెతుకుతూనే ఉండి $S$ను అర్ధ-నిర్ణయించే ట్యూరింగ్
యంత్రం~$G$ను తీసుకోండి.}{}

చివరగా \olref{case:ce-domain} నుంచి~\olref{case:ce} వస్తుందని చూపడానికి,
$S$ అనేది పాక్షిక గణనీయ ప్రమేయం~$\cfind{e}$ యొక్క నిర్వచన క్షేత్రం
అనుకుందాం, అంటే
\[
S=\Setabs{x}{\cfind{e}(x)\fdefined}.
\]
$S$ శూన్య సమితి అయితే పని పూర్తయింది; లేకపోతే $S$లో ఏదైనా మూలకం
$a$ను తీసుకోండి. $f$ను ఇలా నిర్వచించండి:
\[
f(z)=\begin{cases}
(z)_0 & \text{$T(e,(z)_0,(z)_1)$ అయితే} \\
a & \text{ఇతర సందర్భాల్లో.}
\end{cases}
\]
అప్పుడు పైన చెప్పినట్లే, $x$ అనే సంఖ్య $f$ వ్యాప్తిలో ఉంటేనే
$\cfind{e}(x)\fdefined$ అవుతుంది, అంటే $x\in S$ అయితేనే అది $f$
వ్యాప్తిలో ఉంటుంది. \iftag{TMs}{ట్యూరింగ్ యంత్రాల పరంగా: $S$ను
అర్ధ-నిర్ణయించే యంత్రం $M_e$ ఇచ్చినప్పుడు, సాధ్యమైన ట్యూరింగ్ యంత్ర
గణనలన్నిటిలో పరిగెత్తుతూ, ఆగే గణనలకు సంబంధించిన ఇన్‌పుట్‌లను తిరిగి
ఇవ్వడం ద్వారా $S$లోని మూలకాలను లెక్కించండి.}{}
\end{proof}

ఉపవాక్యం~\olref{case:ce-domain} ద్వారా \olref{thm:ce-equiv}, గణనీయంగా
లెక్కించదగిన సమితులను లెక్కించడానికి అనుకూలమైన విధానాన్ని ఇస్తుంది:
ప్రతి~$e$కు, $W_e$ అనేది $\cfind{e}$ యొక్క నిర్వచన క్షేత్రం అని,
అంటే
\[
W_e=\Setabs{x}{\cfind{e}(x)\fdefined}.
\]
అని ఉంచండి. అప్పుడు $A$ ఏదైనా గణనీయంగా లెక్కించదగిన సమితి అయితే,
ఏదో~$e$కు $A=W_e$ అవుతుంది.

గణనీయంగా లెక్కించదగిన సమితులకు మరొక లక్షణీకరణను కింది సిద్ధాంతం
ఇస్తుంది.

\begin{thm}
\ollabel{thm:exists-char}
కింది సమీకరణ నెరవేరే గణనీయ సంబంధం $R(x,y)$ ఒకటి ఉంటే, అప్పుడే మాత్రమే
$S$ గణనీయంగా లెక్కించదగిన సమితి:
\[
S=\Setabs{x}{\lexists[y][R(x,y)]}.
\]
\end{thm}

\begin{proof}
ముందు దిశలో, $S$ గణనీయంగా లెక్కించదగినది అనుకుందాం. అప్పుడు ఏదో
$e$కు $S=W_e$. ఈ~$e$కు $S$ను ఇలా రాయవచ్చు:
\[
S=\Setabs{x}{\lexists[y][T(e,x,y)]}.
\]
విలోమ దిశలో, $S=\Setabs{x}{\lexists[y][R(x,y)]}$ అనుకుందాం.
$f$ను ఇలా నిర్వచించండి:
\[
f(x)\simeq\umin{y}{R(x,y)}.
\]
అప్పుడు $f$ పాక్షిక గణనీయ ప్రమేయం; $S$ దాని నిర్వచన క్షేత్రం.
\end{proof}


\end{document}
